\clearpage\thispagestyle{empty}
\pdfbookmark[0]{Quellen zur Tafel zur Mittagskorrektur (1811)}{mittags1811-sources}
\begin{center}{\textbf{Quellen zur Tafel zur Mittagskorrektur (1811)}}\end{center}
\noindent Diese KI-Übertragung in heutiges Deutsch folgt Carl Friedrich Gauß' deutschem Text in den \href{https://gdz.sub.uni-goettingen.de/id/PPN236007467}{\textit{Werken}, Band~VI} (Göttingen, 1874), S.~166--171. Sie enthält den Brief vom März 1811, das Rechenbeispiel und die drei Tafelseiten der Werkausgabe.
\par\smallskip\noindent Der \href{https://www.digitale-sammlungen.de/en/view/bsb10538612?page=411}{Erstdruck} erschien als Beitrag XXXIV, \textit{Auszug aus einem Schreiben des Herrn Professor Gauß}, im Aprilheft 1811 der von Franz Xaver von Zach herausgegebenen \textit{Monatlichen Correspondenz zur Beförderung der Erd- und Himmels-Kunde}, Band~23, S.~401--409. Die \href{https://www.digitale-sammlungen.de/en/view/bsb10538612?page=414}{Tafel auf S.~404--409} nennt Christian Ludwig Gerling als den Bearbeiter der Berechnungen. Der Erstdruck datiert das Beispiel auf den 2.~März. In der Werkausgabe steht der 20.~März.
\par\smallskip\noindent Die im Brief erwähnten Tafeln von Zachs stehen im \href{https://www.digitale-sammlungen.de/en/view/bsb10538612?page=65}{Januarheft desselben Bandes}, nach S.~56.
\par\smallskip\noindent Übertragung und Satz wurden mit KI erstellt.
\clearpage
